\documentclass[a4paper,twoside]{article}

\usepackage{morewrites}

\usepackage[svgnames,dvipsnames,x11names]{xcolor}
\usepackage{graphicx}
\usepackage[english]{babel}
\usepackage[colorlinks=true, linkcolor=NavyBlue, urlcolor=NavyBlue, citecolor=NavyBlue]{hyperref}
\usepackage[a4paper]{geometry}
\usepackage{ts-fonts}
\usepackage{amsmath}
\usepackage{float}
\usepackage{seqsplit}
\usepackage{ts-code}

\usepackage{lastpage}
\usepackage{refcount}
\makeatletter
\def\ps@plain{
  \let\@oddhead\@empty
  \let\@evenhead\@empty
  \def\@oddfoot{
    \hfil
    \ifnum\getpagerefnumber{LastPage}>1\relax
      \thepage
    \fi
    \hfil}
  \let\@evenfoot\@oddfoot
}
\makeatother

\title{Template Discovery}
\author{}\date{}

\begin{document}

\maketitle

\thispagestyle{plain}
TeXSmith finds templates from multiple locations in a deterministic order:

\begin{enumerate}
\item \textbf{Built-ins}: shipped with TeXSmith (\texttt{article}, \texttt{book}, \texttt{letter}, \texttt{snippet}).
\item \textbf{Installed packages}: PyPI distributions named \texttt{texsmith-\allowbreak{}template-\allowbreak{}*} (or exposing the \texttt{texsmith.templates} entry point).
\item \textbf{Local tree}: current working directory and any ancestor \texttt{templates/} folder. Any \texttt{manifest.toml}/\texttt{template/manifest.toml} or \texttt{\_\_init\_\_.py} counts as a template root.
\item \textbf{User directory}: \texttt{\textasciitilde{}/.texsmith/templates/<name>} (same structure as local).

\end{enumerate}
Use the CLI to inspect what was found:

\begin{code}{bash}{}{}
\begin{Verbatim}[commandchars=\\\{\},breaklines, breakanywhere, commandchars=\\\{\}]
texsmith\PY{+w}{ }\PYZhy{}\PYZhy{}template\PYZhy{}info\PY{+w}{ }\PYZhy{}\PYZhy{}template\PY{+w}{ }article
texsmith\PY{+w}{ }templates\PY{+w}{  }\PY{c+c1}{\PYZsh{} list all visible templates}
\end{Verbatim}
\end{code}
A valid template root contains either \texttt{manifest.toml} or \texttt{template/manifest.toml}; an \texttt{\_\_init\_\_.py} alongside these allows specialized Python logic.

Notes:
- Passing an explicit path (\texttt{-\allowbreak{}-\allowbreak{}template ./templates/custom}) bypasses discovery order.
- Package roots win over same-named local folders; local folders win over the home directory.
- Template manifests can include a \texttt{mermaid-\allowbreak{}config.json} at the root; \texttt{-\allowbreak{}-\allowbreak{}template-\allowbreak{}info} will surface it.

To scaffold a built-in for customization:

\begin{code}{bash}{}{}
\begin{Verbatim}[commandchars=\\\{\},breaklines, breakanywhere, commandchars=\\\{\}]
texsmith\PY{+w}{ }templates\PY{+w}{ }scaffold\PY{+w}{ }article\PY{+w}{ }./templates/article
\end{Verbatim}
\end{code}
Then point \texttt{-\allowbreak{}-\allowbreak{}template} to that path. Any \texttt{mermaid-\allowbreak{}config.json} placed at the template root will be picked up automatically.

\end{document}
